\subsection{Đọc bảng so sánh: mức cải thiện và số lỗi}
\textbf{Thứ hạng và ý nghĩa của hai chỉ số.} Trong Bảng~\ref{tab:results},
AlexNet đứng đầu ở cả accuracy và macro-F1; sau đó là ResNet18, VGG11-BN và
ViT tiny. Cả bốn mạng đều cao hơn sáu baseline. HOG--SVM là baseline mạnh nhất,
đạt 53,19\% accuracy, nên được chọn làm mốc phân tích mức cải thiện trong
Bảng~\ref{tab:gains}. Đây là mốc đối chiếu kết quả, không phải cấu hình được
chọn để điều chỉnh các mạng bằng test.

\begin{table}[H]
\centering\small
\caption{Mức cải thiện so với HOG--SVM trên cùng 10.000 ảnh test; ``đpt'' là điểm phần trăm.}
\label{tab:gains}
\begin{tabular}{lrrrr}
\toprule
Mạng & Tăng accuracy (đpt) & Tăng macro-F1 (đpt) & Số lỗi & Giảm lỗi tương đối \\
\midrule
AlexNet & 30.68 & 30.76 & 1,613 & 65.54\% \\
ResNet18 & 28.26 & 28.37 & 1,855 & 60.37\% \\
VGG11-BN & 26.33 & 26.39 & 2,048 & 56.25\% \\
ViT tiny & 15.02 & 14.98 & 3,179 & 32.09\%
 \\
\bottomrule
\end{tabular}
\end{table}

HOG--SVM phân loại sai 4.681 ảnh; AlexNet sai 1.613 ảnh, tức ít hơn 3.068 ảnh.
Mức tăng accuracy 30,68 điểm phần trăm tương ứng giảm khoảng 65,54\% số lỗi
so với baseline, không phải tăng 30,68\% tương đối về accuracy.
ResNet18, VGG11-BN và ViT tiny lần lượt giảm 2.826, 2.633 và 1.502 lỗi.
Vì cùng dùng tập test 10.000 ảnh, số ảnh đúng và sai có thể đối chiếu trực tiếp.
Tuy nhiên, số liệu tổng hợp không cho biết chính xác ảnh nào được sửa đúng
hay ảnh nào mới bị sai; cần dự đoán ghép theo cùng test index để phân tích đó.

Accuracy và macro-F1 của mỗi mạng khá gần nhau: chênh khoảng 0,06 điểm phần
trăm ở AlexNet, 0,04 ở ResNet18, 0,08 ở VGG11-BN và 0,18 ở ViT tiny.
Cải thiện vì vậy vẫn thể hiện khi mỗi lớp được cho trọng số ngang nhau trong
macro-F1. Điều này không có nghĩa mọi lớp có chất lượng giống nhau; chẳng hạn
AlexNet vẫn có F1 của \emph{cat} thấp hơn rõ rệt so với \emph{automobile}.
SIFT--Random Forest có accuracy 18,57\% nhưng macro-F1 chỉ 14,33\%, gợi ý
hiệu quả phân lớp mất cân đối hơn và cần kiểm tra chỉ số từng lớp.

\subsection{Vì sao học sâu đạt kết quả tốt hơn các baseline?}
\textbf{1. Đặc trưng được học theo mục tiêu phân lớp.} HOG và SIFT cố định
cách lấy gradient hoặc descriptor trước khi học classifier. Nếu thông tin
hữu ích bị bỏ ở bước biểu diễn, SVM hay Random Forest không thể khôi phục từ
vector đã mất thông tin. Với mạng sâu, sai số từ nhãn truyền ngược tới các
lớp đầu, điều chỉnh cả bộ lọc lẫn bộ phân lớp. Mạng có thể chọn những tổ hợp
cạnh, màu và kết cấu phù hợp với 10 lớp trong train, thay vì dùng cùng một
quy tắc trích xuất thủ công cho mọi đối tượng. Trong bảng, ưu thế xuất hiện
ở cả ba CNN lẫn ViT, phù hợp với lợi ích của học biểu diễn theo nhiệm vụ.

\textbf{2. Khai thác không gian và thông tin màu.} Các mạng nhận ảnh RGB và
xử lý cấu trúc ảnh thông qua tích chập hoặc embedding vị trí patch.
HOG và SIFT ở đây dùng ảnh xám, nên không trực tiếp giữ khác biệt giữa các
mảng màu có mức sáng tương tự. Pixel thô vẫn giữ RGB và thứ tự vị trí,
nhưng Linear SVM chỉ áp dụng một tổ hợp tuyến tính trên toàn vector.
Khi vị trí vật thể thay đổi, một bộ trọng số gắn với từng tọa độ cố định có
thể kém linh hoạt. CNN tái sử dụng cùng bộ lọc ở nhiều vị trí; pooling và
các tầng sau kết hợp đáp ứng cục bộ. ViT giữ vị trí token và dùng attention
để kết hợp các vùng. Ưu thế này không đồng nghĩa mạng bất biến hoàn toàn với
mọi phép dịch chuyển, xoay hay thay đổi nền.

\textbf{3. Tạo tổ hợp phi tuyến nhiều tầng.} SVM tuyến tính không tự tạo
các tương tác phi tuyến mới giữa chiều đặc trưng. Random Forest có thể học
quy tắc phi tuyến, nhưng các phép chia một chiều không cung cấp cấu trúc học
mẫu không gian như convolution hay attention. Ở mạng sâu, convolution/linear,
ReLU hoặc GELU và các tầng liên tiếp tạo nhiều mức tổ hợp. Chẳng hạn, một
đường viền đơn lẻ có thể xuất hiện ở nhiều lớp; kết hợp đường viền với màu,
vị trí và các vùng lân cận có thể giúp phân biệt tốt hơn. Điều này giải thích
về mặt cơ chế vì sao tăng độ linh hoạt của biểu diễn có thể nâng accuracy.

\textbf{4. Tránh vấn đề thứ tự descriptor của SIFT phẳng.} SIFT--SVM đạt
22,92\% và SIFT--Random Forest đạt 18,57\%, thấp hơn cả các baseline pixel.
Trong vector SIFT nối trực tiếp, descriptor thứ nhất của một ảnh có thể nằm
ở cánh máy bay, nhưng descriptor thứ nhất của ảnh khác nằm ở nền. Việc đệm
0 còn tạo khác biệt theo số keypoint. Trái lại, CNN giữ lưới không gian qua
các feature map; ViT gán embedding vị trí và có phép kết hợp token học được.
Cả hai không yêu cầu các điểm mạnh nhất phải xuất hiện theo cùng thứ tự.
Kết quả thấp này là hạn chế của SIFT phẳng rút gọn đang dùng, không phải bằng
chứng rằng mọi pipeline SIFT đều kém; SIFT--BoVW hoặc spatial pooling có thể
thay đổi biểu diễn và cần được đánh giá riêng.

\textbf{5. Số tham số, tối ưu và điều chuẩn cũng khác nhau.} Các mạng có
khoảng 2,31--2,80 triệu tham số và được tối ưu bằng AdamW, trong khi SVM là
mô hình tuyến tính và Random Forest dùng cấu trúc cây. Chuẩn hóa activation,
weight decay, gradient clipping và dropout trong AlexNet góp phần định hình
quá trình học. Do đó, bảng chứng minh hiệu quả của toàn bộ cấu hình học sâu
đã chạy; chưa tách riêng ảnh hưởng của học đặc trưng, giữ RGB, kích thước mô
hình hay optimizer. Muốn xác định nguyên nhân định lượng, cần ablation như
CNN trên ảnh xám, classifier tuyến tính trên embedding CNN cố định, hoặc
thay đổi một thành phần điều chuẩn trong khi giữ các yếu tố khác giống nhau.

\subsection{Vì sao thứ hạng giữa các mô hình lại như trong bảng?}
\textbf{Trong nhóm baseline.} Pixel--Random Forest đạt 45,09\%, hơn
Pixel--SVM 7,52 điểm phần trăm, phù hợp với lợi ích của ranh giới phi tuyến
trên đầu vào màu thô. Khi chuyển sang HOG, SVM đạt 53,19\%, cao hơn Random
Forest 6,46 điểm phần trăm. Đặc trưng histogram nhiều chiều đã mã hóa hình
dạng ở lưới cố định, có thể phù hợp hơn với bộ phân lớp dùng đồng thời nhiều
chiều. Random Forest giới hạn độ sâu 15 và xét một tập chiều nhỏ ở mỗi nút;
việc mô hình hóa các tổ hợp rộng của gradient có thể đòi hỏi nhiều phép chia.
Các lý giải này cần kiểm tra bằng điều chỉnh siêu tham số, không đủ để khẳng
định SVM luôn tốt hơn Random Forest trên HOG.

\textbf{AlexNet và hai CNN còn lại.} AlexNet cao hơn ResNet18 2,42 điểm phần
trăm và VGG11-BN 4,35 điểm phần trăm accuracy. AlexNet có GroupNorm, dropout
0,5 ở classifier và còn giữ lưới $2\times2$ sau adaptive pooling; VGG dùng
pooling tới $1\times1$, ResNet dùng global average pooling. Các lựa chọn
này ảnh hưởng cách giữ vị trí, khả năng biểu diễn và điều chuẩn. VGG có ít
tham số nhất, còn ResNet18 có nhiều nhất, nhưng số tham số không xếp đúng
thứ hạng accuracy: lớn hơn hoặc sâu hơn không bảo đảm tốt hơn.
ResNet giúp tối ưu mạng qua shortcut; ưu điểm đó không tự giải quyết được
quá khớp hay bảo đảm lịch tối ưu chung là phù hợp nhất cho ResNet.
Chưa có ablation hoặc nhiều seed để xác định yếu tố nào quyết định mức chênh.

\textbf{ViT tiny so với CNN.} ViT có số tham số gần AlexNet nhưng accuracy
thấp hơn 15,66 điểm phần trăm. CNN có thiên kiến cấu trúc về tính cục bộ và
chia sẻ bộ lọc; ViT dùng tương tác toàn chuỗi patch và phải học nhiều quan hệ
từ dữ liệu. Nghiên cứu ViT chỉ ra vai trò của lượng dữ liệu và thiên kiến cấu
trúc khi huấn luyện \cite{dosovitskiy2021vit}. Trong thí nghiệm này chỉ có
45.000 ảnh train, không augmentation, không pretrained; ViT cũng không có
dropout trong encoder. Các điều kiện này phù hợp với giả thuyết rằng ViT
khó tổng quát hóa hơn ba CNN, dù loss train giảm rất thấp. Attention cho phép
nhìn toàn ảnh nhưng không bảo đảm sử dụng đúng dấu hiệu vật thể; nền vẫn có
thể trở thành tín hiệu gây nhầm. Kết quả 68,21\% vẫn cao hơn HOG--SVM, nên
ViT học được biểu diễn hữu ích trong cấu hình này.

\textbf{Hiệu quả và giới hạn phép so sánh.} VGG11-BN mất khoảng 2,19 giờ,
ngắn hơn AlexNet 1,29 giờ, nhưng giảm 4,35 điểm phần trăm accuracy.
ResNet18 mất khoảng 6,51 giờ, dài hơn AlexNet khoảng 3,03 giờ, trong khi
accuracy thấp hơn. Vì vậy, AlexNet có sự cân bằng thuận lợi giữa accuracy
và thời gian trong lần chạy này. Thời gian train không thay thế thời gian
inference, bộ nhớ cực đại hay chi phí năng lượng. Các mạng dùng MPS; baseline
NumPy có cài đặt và thiết bị thực thi khác. Cùng số mẫu và seed giúp đối chiếu
dữ liệu, còn 200 epoch và 200 cây không tạo ngân sách tính toán ngang nhau.
Một lần chạy chưa đủ để đánh giá độ ổn định hay ý nghĩa thống kê của thứ hạng.
